\section{来源审计：两种把语言接到机器人上的接口}

这堂课讨论的不是泛泛的“机器人接入大模型”，而是一个更具体的 systems question：高层语言和视觉语义已经很强，怎样把它们可靠地接到连续、实时、受物理约束的低层动作？Fei Xia 用两条路线回答。第一条是 PaLM-E、RT-2 与 RT-X 所代表的\textbf{模型整合路线}：把 perception、language reasoning 与 action prediction 收进一个可共同训练的模型。第二条是 Language to Rewards 所代表的\textbf{接口分层路线}：让 LLM 生成 reward program，再由 Model Predictive Control（MPC，模型预测控制）或其他优化器执行低层控制。

\subsection{官方课堂边界与旧稿问题}

本次重写以 Stanford Online 官方录像 \texttt{fz8wf9hN20c} 为唯一课堂顺序来源，课堂日期为 2023 年 10 月 10 日。录像提供人工 \texttt{en-US} 字幕，55 个完整教学状态从 1080p 视频逐页恢复。旧稿没有视频链接，只有一张正文图，并把 PaLM-E、RT-2、reward generation 和安全问题压缩成产品名词列表；它没有呈现 lecture 的两段结构、数据/接口两个根因、action token 的机制、co-fine-tuning 的证据链，也没有保留问答中的限制条件，因此本次直接重建而非局部扩写。

\begin{longtable}{@{}L{0.18\textwidth}L{0.31\textwidth}L{0.39\textwidth}@{}}
\toprule
证据层 & 本地材料 & 在讲义中的职责 \\
\midrule
官方录像 & 1:18:13、55 个最终教学状态 & 决定 slide 顺序、demo、图表、公式和 Q\&A 边界 \\
官方人工字幕 & 1,597 条 caption、timed transcript & 恢复动机、警告、设计取舍、安全与未来方向 \\
Primary sources & PaLM-E、RT-1、RT-2、RT-X、Language Table、Language to Rewards & 核对模型结构、训练目标、benchmark 与术语 \\
视觉审计 & 55 张全宽截图与 selection TSV & 同一视频 slide 只保留代表状态，不把逐帧播放当成不同页面 \\
\bottomrule
\end{longtable}

\begin{warningbox}{时间边界与证据边界}
本文重建 2023-10-10 的课堂，不把后续机器人模型倒灌为课堂事实。视频 demo 说明某个系统在特定设置中完成过任务，不等于可靠处理开放世界；benchmark 提升说明相对基线更好，不等于具备人类级物理常识；“positive transfer”说明联合训练在给定任务混合中带来收益，也不等于跨所有机器人、任务与环境都能迁移。
\end{warningbox}

\subsection{读图路线：从物理失败到接口选择}

标题页把 Foundation Models 放在 Low-level Embodied Intelligence 的中心，但真正的主线不是“模型越大越好”。读图时应持续追踪五个对象：\textbf{observation} 提供了什么、\textbf{model} 学到了什么、\textbf{interface} 输出何种表示、\textbf{controller} 如何把表示变成动作、\textbf{evidence} 在什么分布上成立。后面所有架构和结果都可以放回这五格检查表。

\lecturefigure{slide-01-title-low-level-embodied.jpg}{Lecture 19：Foundation Models 与低层具身智能}{官方录像恢复幻灯片 V001；00:01:18}

\begin{importantbox}{本讲的核心判断}
语言模型不是直接替代机器人系统，而是改变系统中\textbf{语义知识的来源}与\textbf{任务接口的表达方式}。低层控制仍要面对动作频率、动力学、接触、约束、延迟、硬件极限和安全层；真正的工程问题是把 foundation-model prior 接到这些物理环节，而不是忽略它们。
\end{importantbox}

\subsection{本章小结}

本章确定了课堂边界与阅读坐标。后文先解释物理世界为何难，再介绍模型整合路线 PaLM-E/RT-2/RT-X，随后转向 Language to Rewards，最后用 Q\&A 比较两种接口的收益、失败模式与混合方案。

